\originalsection{第8次习题课}
\sourceinfo{MIT 18.04, Spring 2018, Recitation 8; 题面 \nolinkurl{mit18_04s18_recit8-handout.pdf}, PDF第1页; 解答 \nolinkurl{mit18_04s18_recit8-solutions.pdf}, PDF第1页. 课程作者 Jeremy Orloff, 页内署名 Vishesh Jain, 2018, CC BY-NC-SA 4.0}
\begin{exercise}
\textbf{原题 Rec8.1。}\label{ass:rec8-1}
1.1. 若 $g$ 在 $z_0$ 有简单零点, 证明 $1/g$ 在 $z_0$ 有简单极点. 1.2. 证明 $\Res(1/g,z_0)=1/g'(z_0)$. 1.3. 找出 $1/\sin z$ 的全部极点, 证明均为简单极点并求留数.
\end{exercise}
\begin{solution}
编者补充 (AI): 官方解答仅指向 Topic 8 第5页性质5与例8.11. 简单零点意味着 $g(z)=(z-z_0)h(z)$, 其中 $h$ 在邻域解析且 $h(z_0)=g'(z_0)\ne0$. 因而 $1/g=(z-z_0)^{-1}/h(z)$ 有简单极点, 留数为 $\lim_{z\to z_0}(z-z_0)/g(z)=1/g'(z_0)$. $\sin z=0$ 等价于 $\ee^{2\ii z}=1$, 其全部解为 $z=k\pi$, $k\in\Z$. 在这些点 $\cos(k\pi)=(-1)^k\ne0$, 所以零点均简单, $1/\sin z$ 的全部极点恰为这些点, 留数均为 $(-1)^k$.
\end{solution}
\begin{exercise}
\textbf{原题 Rec8.2。}\label{ass:rec8-2}
2.1. 设 $p,q$ 在 $z_0$ 解析, $p(z_0)\ne0$, $q$ 在 $z_0$ 有简单零点. 证明 $\Res(p/q,z_0)=p(z_0)/q'(z_0)$. 2.2. 找出 $\cot z$ 的全部极点, 证明均为简单极点并求留数.
\end{exercise}
\begin{solution}
编者补充 (AI): 官方解答仅指向例8.13与8.4.3节. 写 $q(z)=(z-z_0)h(z)$, $h(z_0)=q'(z_0)\ne0$. 由于 $p/h$ 解析且在 $z_0$ 非零, $p/q$ 的简单极点留数为 $p(z_0)/h(z_0)$. 对 $\cot z=\cos z/\sin z$, 分母的简单零点为 $k\pi$, 分子在每一点均为 $(-1)^k\ne0$, 所以全部极点为 $k\pi$, 每个留数均为 $\cos(k\pi)/\cos(k\pi)=1$.
\end{solution}
\begin{exercise}
\textbf{原题 Rec8.3。}\label{ass:rec8-3}
利用 $\cos z,\sin z$ 在0处的 Taylor 级数, 求 $\cot z$ 在0处 Laurent 展开的前几项.
\end{exercise}
\begin{solution}
编者补充 (AI): 官方解答仅指向例8.17. 因 $\cot z$ 为奇函数且留数为1, 写 $\cot z=z^{-1}+az+bz^3+cz^5+O(z^7)$. 与 $\sin z=z-z^3/6+z^5/120-z^7/5040+O(z^9)$ 相乘, 和 $\cos z=1-z^2/2+z^4/24-z^6/720+O(z^8)$ 比较, 得
\[
a-\tfrac16=-\tfrac12,\qquad b-\tfrac a6+\tfrac1{120}=\tfrac1{24},\qquad
c-\tfrac b6+\tfrac a{120}-\tfrac1{5040}=-\tfrac1{720}.
\]
故 $a=-1/3$, $b=-1/45$, $c=-2/945$, 即 $\cot z=z^{-1}-z/3-z^3/45-2z^5/945+O(z^7)$. 最近的其他极点为 $\pm\pi$, 所以该 Laurent 级数在 $0<|z|<\pi$ 收敛.
\end{solution}
\begin{exercise}
\textbf{原题 Rec8.4。}\label{ass:rec8-4}
设 $f$ 在区域 $A$ 中除一组孤立奇点外解析, $C$ 是 $A$ 内不经过奇点的逆时针简单闭曲线. 4.1. 若 $C$ 内只有奇点 $z_1$, 用扩展的 Cauchy 定理证明 $\int_Cf(z)\dd z=2\pi\ii\Res(f,z_1)$. 4.2. 若内部有两个奇点 $z_1,z_2$, 证明积分为 $2\pi\ii(\Res(f,z_1)+\Res(f,z_2))$. 4.3. 推广为内部所有留数之和的 $2\pi\ii$ 倍, 即留数定理.
\end{exercise}
\begin{solution}
编者补充 (AI): 官方解答PDF保留4.1--4.3题面却没有给出答案. 原题还须补充 \textbf{$C$ 及其所围闭区域均包含在 $A$ 内, 且曲线分段光滑}, 以下在该条件下证明. 只要求 $C\subset A$ 不够: 例如在环域 $A=\{1<|z|<3\}$ 上 $f=1/z$ 解析, 沿 $|z|=2$ 的积分却为 $2\pi\ii$, 不能由 $A$ 中的奇点求和得到0.

4.1取以 $z_1$ 为圆心的小圆 $c_1$, 使闭圆盘位于 $C$ 内且不含其他奇点, $c_1$ 本身逆时针. 在外区域挖去该小圆盘, $f$ 在所得多连通区域及边界解析. 正向边界为 $C-c_1$, 扩展 Cauchy 定理给 $\int_C f=\int_{c_1}f$. 若 $f(z)=\sum_{k\in\Z}a_k(z-z_1)^k$ 是小圆附近的 Laurent 展开, 其在小圆上一致收敛, 可逐项积分. 除 $k=-1$ 外各项积分为0, 所以 $\int_{c_1}f=2\pi\ii a_{-1}=2\pi\ii\Res(f,z_1)$.

4.2分别挖去互不相交的小圆盘, 正向边界为 $C-c_1-c_2$. 扩展 Cauchy 定理给 $\int_Cf=\int_{c_1}f+\int_{c_2}f$, 再对两个小圆用同样的 Laurent 计算, 得到题给公式.

4.3闭内部是 $A$ 的紧子集, 孤立奇点在 $A$ 内无聚点, 所以内部只能有有限多个奇点 $z_1,\ldots,z_m$. 取互不相交且包含在内部的小圆盘, 挖去后正向边界是 $C-\sum_jc_j$. 因而 $\int_Cf=\sum_j\int_{c_j}f=2\pi\ii\sum_{j=1}^m\Res(f,z_j)$. 求和为空时即为通常的 Cauchy 定理.
\end{solution}
